% Image credits for the photographs and AI illustrations of Book 4
% (University Chemistry, Year 3), Dutch edition. Machine-readable license
% records live in images/book4/CREDITS.md; keep this file in sync with
% frontmatter/image-credits-book4.tex.
\chapter*{Beeldverantwoording}

De tekeningen in dit boek zijn origineel. Foto's worden, waar ze gebruikt
worden, gebruikt onder hun respectieve vrije licenties, met dank aan hun
makers; de volledige bronlinks en licentie-URL's van elke foto staan in
\texttt{images/book4/CREDITS.md} in de repository van het boek.
De gevarenpictogrammen zijn die van het Globally Harmonized System of
Classification and Labelling of Chemicals, zoals getekend voor het pakket
\texttt{ghsystem} van \TeX\ Live.

\bigskip
Naast de foto's en de originele tekeningen zijn sommige figuren
illustraties die met AI gegenereerd zijn, gemaakt met de beeldgeneratie
van OpenAI op basis van prompts die door de redactie van het boek
geschreven zijn, en vóór opname nagekeken op chemische juistheid. De
volledige prompt van elk gegenereerd beeld staat in
\texttt{images/book4/ai/PROMPTS.md} in de repository.
\bigskip
\noindent Foto's, per hoofdstuk:
\begin{itemize}\small
  \item Hfst.~1: Erwin Schr\"odinger, 1933, Nobel Foundation, publiek domein
        (Wikimedia Commons).
  \item Hfst.~4: sneeuwkristallen, Wilson A. Bentley (rond 1902), publiek
        domein (Wikimedia Commons).
  \item Hfst.~6: radioantennes op een hoogvlakte, ESO/B.~Tafreshi
        (twanight.org), CC BY 4.0 (Wikimedia Commons).
  \item Hfst.~7: fluoriet bij daglicht en onder ultraviolet licht, door Masha
        Milshina, CC BY 4.0 (Wikimedia Commons).
  \item Hfst.~8: een NMR-laboratorium, door Shandchem, CC BY 2.0 (Wikimedia
        Commons).
  \item Hfst.~9: William Lawrence Bragg, Nobel Foundation, publiek domein
        (Wikimedia Commons).
  \item Hfst.~10: het graf van Ludwig Boltzmann, door Feldkurat Katz,
        CC BY-SA 4.0 (Wikimedia Commons).
  \item Hfst.~13: spiralen van Belousov-Zhabotinsky, door Stephen Morris,
        CC BY 2.0 (Wikimedia Commons).
  \item Hfst.~14: het ozongat boven Antarctica, NASA, publiek domein
        (Wikimedia Commons).
  \item Hfst.~16: elektronenmicroscopische opname van de honingraat van een
        katalysator, door Galadrid, CC BY-SA 4.0 (Wikimedia Commons).
  \item Hfst.~17: zonnestralen en het tyndalleffect, door Kevin c higgins,
        CC BY-SA 4.0 (Wikimedia Commons).
  \item Hfst.~18: robijn- en smaragdkristallen, door Robert M. Lavinsky,
        CC BY-SA 3.0 (Wikimedia Commons).
  \item Hfst.~20: kristallen van ferroceen, door Andrea Sella, CC BY 4.0
        (Wikimedia Commons).
  \item Hfst.~22: een éénkristal van silicium, door Sebastian Wallroth,
        publiek domein (Wikimedia Commons).
  \item Hfst.~23: een bloem op silica-aerogel, NASA/JPL, publiek domein; de
        Lycurgusbeker, door Marie-Lan Nguyen, CC BY 2.5 (Wikimedia Commons).
  \item Hfst.~24: een Atlantische degenkrab, door Kaldari, CC0 (Wikimedia
        Commons).
  \item Hfst.~27: pyrethrummargrieten, door Soramimi, CC BY-SA 4.0
        (Wikimedia Commons).
  \item Hfst.~30: schors van een westelijke taxus, door JOE BLOWE
        (theforestprimeval), CC BY-SA 2.0 (Wikimedia Commons).
  \item Hfst.~32: algenbloei gezien door de satelliet Landsat 8, USGS en
        NASA, publiek domein.
  \item Hfst.~33: een dubbel verdeelstuk voor vacuüm en inert gas, door
        Polimerek, CC BY-SA 3.0 (Wikimedia Commons).
\end{itemize}
\clearpage
